\subsection{Optimizing among atheoretical \persona{s}}
\label{sec:atheory_prompt}

Similarly to our study of AR, grounding a \persona{} template in theory is important for generalization.
We do this via negative example.
Specifically, without a theoretical grounding, the optimization procedure may fail to find a \persona{} that even fits in-sample.

Unlike in the Section~\ref{sec:predict_new_AR} with the Always Pick `N' agents, we do not offer an analogous overfitting example.
Because human data from multiple games is used for optimization, finding a sample of \aisub{s} which overfits requires finding \persona{s} which overfit to all six settings.
This is a much more difficult task than finding a \persona{} which overfits to a single game.
Indeed, this is an attractive feature of using multiple training samples to construct and validate agent samples.

To generate samples of arbitrary agents, we repeat the entire process from Section~\ref{sec:apply_prompt} but replace the theory-grounded attributes (i.e., efficiency, inequity aversion, and self-interest) with wholly unscientific ones: a self-reported fondness for the TV show \emph{New Girl}, an enthusiasm for taxidermy, and swimming ability.  
This new prompt template is:
\begin{quote}
\vspace{-.5cm}
\singlespacing
$\theta(\phi_{ng}, \phi_{tax}, \phi_{swim})$ = \textit{On a scale from 1 to 10, you think the show New Girl is: \{$\phi_{ng}$\}.
10 means you love New Girl, and 1 means you hate it.
On a scale from 1 to 10, your passion for taxidermy is: \{\(\phi_{tax}\)\}.
10 means you love taxidermy, and 1 means you hate it.
On a scale from 1 to 10, your ability to swim is: \{\(\phi_{swim}\)\}.
10 means you are a great swimmer, and 1 means you can't swim.}
\end{quote}

Using identical hyperparameters and the same Bayesian optimization procedure, we search over this atheoretical space to see if any combination of $(\phi_{ng}, \phi_{tax}, \phi_{swim})$ (each a sample of 3 agents with their own parameter vector) could even match the original single-stage dictator games in-sample.  
The resulting atheoretical parameter vector was: $\bigl(\bm{\phi}_{ath-1}^*,\bm{\phi}_{ath-2}^*,\bm{\phi}_{ath-3}^*\bigr) \;=\; \bigl(\,(5,\,7,\,1),\;(9,\,9,\,5),\;(7,\,6,\,8)\bigr)$.
As shown in Figure~\ref{fig:cr_train} (yellow), $\bm{\theta}_{ath}^*$ constructed using these parameters and template failed to beat even the baseline \aisub{s}' performance.  
In fact, throughout the search, no parameter combination for ``loving \emph{New Girl},'' ``passion for taxidermy,'' or ``swimming skill'' ever produced a distribution of choices that aligned more closely with real humans.
This result demonstrates the importance of grounding \aisub{s} in theoretical constructs.

This lack of improvement persisted in the two-stage validation games as well (Figure~\ref{fig:cr_test}; rightmost column).  
The atheoretical \aisub{s} and the baseline were effectively indistinguishable in their distribution of responses as Player A, and the atheoretical subjects were only a little better as Player B.
Overall, the atheoretical \aisub{s} were far less aligned than the theory-grounded \aisub{s}. 
They were closer to the human data than the baseline in $\CRTestAtheoryTotalProp$\% of the settings, worse than the baseline in $\CRTestBaselineTotalPropAtheo$\% of the settings, and identical in the remaining games.
The only way these arbitrary \persona{s} generalize is that their poor performance is consistent across settings.
